% Migration IDs: P-07-001 through P-07-012 and A-07-* in the manifest.
\section{Registered computational study}
\label{sec:computational-study}
\label{sec:experiments}

The theory identifies what each algorithm guarantees, but it does not say
which structural instances are computationally difficult or how much burden
is lost by a heuristic in a finite design. Registered Algorithmic Compression
Benchmark v2 addresses four questions: when do the independent exact methods
agree; how do runtime and diagnostic counts vary with problem size and
structure; how often do certified heuristics attain the available reference
burden; and how much model reduction and endpoint improvement come from
preprocessing and multi-start deletion? The study is synthetic and
identity-closure only. It makes no population claim about real financial
libraries.

\subsection{Registered design and execution}

The final registry contains \AORBenchmarkInstances\ analysis instances:
\AORBenchmarkSmallExact\ small exact instances,
\AORBenchmarkStructured\ structured finite-design instances, and
\AORBenchmarkAdversarial\ adversarial mechanism instances. Family A is
small enough for the registered exact comparisons. Family B varies strategy
and tagged-row counts, coverage density, module overlap, bundle prevalence,
unique-carrier frequency, weight dispersion, and frontier--module coverage
correlation. Family C contains the predeclared gap, tie, duplicate,
dominance, rare-carrier, symmetry, and bundle-versus-singleton mechanisms.
Realized rather than requested structural statistics enter the analysis.

The algorithm grid comprises complete enumeration, requirement-mask dynamic
programming, both registered HiGHS formulations, weighted and cardinality
greedy construction, greedy followed by reverse deletion, five deterministic
rechecked-deletion orders, a seeded random order, and multi-start random
deletion. Mandatory-only and full fixed-point preprocessing variants are
retained as registered. Every returned library is reconstructed in original
coordinates and rechecked for the inactive strategy, exact tagged coverage,
frontier equality, identity closure, and exact burden.

Execution used Julia \AORJuliaVersion, HiGHS \AORHiGHSVersion, and an
\AORCPU. The locked supervisor dispatched \AORWorkerCount\ isolated serial
workers; each worker used serial BLAS and serial HiGHS. Thus wall-clock time
is conditional on the registered concurrent environment, including cache,
memory-bandwidth, and CPU contention. The execution record also preserves
that the worktree was dirty at launch, together with its status hash. That
provenance limitation prevents interpreting these times as a pristine-release
benchmark, although the result files and audit hashes remain internally
traceable. HiGHS used presolve, deterministic registered seeds, and zero
requested relative and absolute MIP gaps. A solver status is still solver
evidence, not an exhaustive or formal proof.

\subsection{Audit and exact-method validation}

The independent audit passed all \AORBenchmarkRuns\ registered run units.
It accepted \AORBenchmarkAccepted\ returned libraries after exact semantic
and burden checks and retained \AORBenchmarkUnsuccessful\ unsuccessful
records. Among \AORExactCapable\ exact-capable instances, complete
cross-method burden agreement was available on \AORExactAgreement; the
remaining \AORExactUnavailable\ had a recorded unavailable method rather
than an imputed optimum. Optimizer identities differed at equal burden on
\AOROptimizerIdentityDifferences\ agreed instances. Such identity ties are
valid alternative optimizers, not disagreements. These statements are exact
finite computation plus independent rechecking; they are not Lean
verification.

Family B has \AORStructuredReferences\ exactly feasible,
HiGHS-\texttt{OPTIMAL} references. Those rows support a consistent finite
comparison but do not upgrade the solver claim. No second solver was pinned.
The complete instance-level agreement and certificate hashes are in Online
Resource~1.

\subsection{Scaling and diagnostic counts}

Table~\ref{tab:aor-computational-scaling-compact} reports the compact scaling
view. Enumeration visits all subsets after mandatory preprocessing, whereas
the DP state count depends on reachable requirement masks. The observed DP
state medians remain far below the worst-case mask bound in these small cells;
this is a finite-design observation and does not change the exponential
complexity proved in Section~\ref{sec:algorithms}. In the structured cells,
median MIP wall time rises when the strategy count moves to the larger level,
while the exposed node medians are small because fixed-point preprocessing
often resolves most or all of the cover before branch-and-bound. Node counts
are unavailable for most such runs and are never replaced by zero.

\AORComputationalScalingTable

The registered penalized-runtime models use log PAR--2, keep every run unit,
and include the declared structural effects, replicate block, fixed lane, and
only the registered algorithm-by-size interactions. The structured model
associates larger strategy count with longer penalized runtime after these
controls. The small-exact and adversarial solved-probability fits encountered
separation, while the structured solved-probability model converged. These
finite-design associations are descriptive; they are not causal performance
laws or cross-machine predictions. Complete coefficients, HC3 intervals, and
Holm-adjusted probabilities are assigned to Online Resource~1.

\subsection{Heuristic quality and structural difficulty}

The two evidence classes in
Table~\ref{tab:aor-heuristic-quality-compact} are intentionally separate.
On Family A, multi-start deletion attains every available exact reference;
heaviest-safe-first and maximum immediate release also perform well in this
random small design. On Family B, weighted greedy followed by reverse deletion
has the highest reference-attainment count among the reported deterministic
constructions and small upper-tail gaps, whereas lightest-safe-first, declared
source order, and a single random deletion order show substantially larger
tails. These are comparative synthetic outcomes, not guarantees.

\AORHeuristicQualityTable

The registered relative-gap models reinforce the need for method-specific
interpretation. Greater unique-carrier frequency is associated with smaller
gaps for several deletion rules, and greater weight dispersion with larger
gaps for cardinality, lightest-first, source-order, random-order, and
multi-start variants. Bundle prevalence is negatively associated with the
weighted-greedy gap. Other declared structural terms are not uniformly
selected after multiplicity adjustment, so the evidence does not support a
single scalar notion of instance difficulty. The separately requested
regression for the probability of reference attainment was not prespecified;
it remains labeled \textsc{exploratory} in Online Resource~1 and is not used
for a confirmatory claim here.

\subsection{Preprocessing, multi-start deletion, and adversarial mechanisms}

Full fixed-point preprocessing has a strong structural effect: the median
variable and requirement reductions in every structured size cell are
complete among available preprocessing records. That does not translate
mechanically into lower end-to-end runtime. The registered full pipeline
includes construction, audit-trail, and reconstruction overhead; its median
log PAR--2 ratio relative to mandatory-only preprocessing is positive in every
structured cell. The two prespecified preprocessing regressions are rank
deficient, so no adjusted causal or universal speedup claim is made.

Multi-start deletion completes all registered starts on
\AORMultistartComplete\ instances and improves over its first start on
\AORMultistartImproved. Improvement is therefore common but not automatic,
and the primary mandatory-only workflow includes the recorded memory-limit
failures rather than substituting successful sensitivity runs.

\AORPreprocessingMultistartTable

The adversarial family connects the experiment back to the theorem without
using experiment to establish it. For the registered heaviest-safe-first gap
instances, the exactly rechecked observed ratios equal the formula from
Theorem~\ref{thm:aor-heaviest-safe-first-gap} at every registered scale. The
small reference is exact finite computation; the larger references remain
HiGHS solver evidence plus exact feasibility checks.

\AORAdversarialGapTable

\subsection{Timeouts, failures, and limitations}

No final run ended at \texttt{TIME\_LIMIT}, so the registered timeout
best-bound-gap table is correctly empty. This is not evidence that time limits
are irrelevant: \AORGenerationFailed\ run units correspond to rejected
structured generations, \AORMemoryLimited\ multi-start run units reached
the memory cap, and implementation, interruption, missing-primal, and
nonbinary-candidate failures remain in the machine-readable census. Solved
fractions and PAR--2 summaries retain those outcomes in their denominators.

The remaining boundaries are substantive. The benchmark covers finite tagged
identity-closure instances, not arbitrary closure. Its structural grid is not
a probability sample of financial libraries. Runtime is machine- and
contention-conditional, MIP reference quality rests on one open-source solver,
and unavailable diagnostics remain missing. Accordingly, the study identifies
where the implemented methods differ within the locked design; it does not
prove average-case complexity or a universal ranking of algorithms.
